Los hitos del proyecto representan los momentos clave dentro del ciclo de desarrollo del software, en los cuales se finaliza o valida un entregable importante que permite avanzar hacia la siguiente etapa del proyecto.

\begin{itemize}

\item \textbf{Acta de inicio del proyecto} \
Se formaliza el inicio del proyecto de desarrollo del sistema, estableciendo los objetivos generales, el alcance preliminar y los actores involucrados.

\item \textbf{Validación de la propuesta del proyecto} \
Se revisa y aprueba la propuesta de desarrollo del sistema por parte de los responsables académicos o interesados, confirmando la viabilidad del proyecto.

\item \textbf{Levantamiento de requerimientos} \
Se identifican, documentan y analizan los requerimientos funcionales y no funcionales que deberá cumplir el sistema para optimizar el proceso de gestión de licencias.

\item \textbf{Especificación de requerimientos del sistema} \
Se consolida el documento de especificación de requerimientos del software, donde se detallan formalmente las funcionalidades, restricciones y características del sistema.

\item \textbf{Diseño de la arquitectura del sistema} \
Se define la arquitectura del software, incluyendo los componentes principales del sistema, la estructura de la aplicación y la interacción entre los módulos.

\item \textbf{Diseño del modelo de datos} \
Se diseña la estructura de la base de datos que soportará el sistema, definiendo entidades, relaciones y reglas de integridad necesarias para la gestión de la información.

\item \textbf{Desarrollo de los módulos del sistema} \
Se implementan las funcionalidades principales del sistema de acuerdo con los requerimientos y el diseño definido.

\item \textbf{Pruebas del sistema} \
Se realizan pruebas funcionales y de integración para verificar que el sistema cumple con los requerimientos establecidos.

\item \textbf{Entrega del sistema desarrollado} \
Se presenta la versión final del sistema desarrollado, incluyendo la documentación técnica y funcional correspondiente.

\end{itemize}
